%% kapa: exam=1 round=31 year=2020 period=1 form=A subject=civil_law subject_ko=민법 seq=38 number=38
%% source: Q-Net 감정평가사 기출문제 (cst00309) artlSeq=5207219 file=extracted/1008/5207219/A형/A형/2020년도 제31회 감정평가사 1차 시험_1교시_A형.pdf
%% conversion: kordoc 4.24.51 -> tools/gen1.py (auto); review=proofread-sonnet; flags=none
\begin{question}{38}
질권에 관한 설명으로 옳지 않은 것은? (다툼이 있으면 판례에 따름)
\begin{choices}
  \item 질권은 질물 전부에 효력이 미친다.
  \item 저당권으로 담보된 채권에 설정된 질권은 그 저당권등기에 질권의 부기등기를 하여야 저당권에 효력이 미친다.
  \item 금전채권에 질권을 취득한 질권자는 자기채권액의 범위에서 직접 추심하여 변제에 충당할 수 있다.
  \item 질권설정자는 피담보채무의 변제기 이후의 약정으로 질권자에게 변제에 갈음하여 질물의 소유권을 이전할 수 있다.
  \item 금전채무자가 채권자에게 담보물을 제공한 경우, 특별한 사정이 없으면 채무자의 변제의무와 채권자의 담보물반환의무는 동시이행관계에 있다.
\end{choices}
\end{question}
\answer{5}
